\section{Introduction}
\label{sec:intro}

Phishing represents a pervasive threat in modern cyberspace. Attackers engineer deceptive URLs and spoofed web portals to impersonate authentic organizations, harvesting credentials and financial assets. Threat telemetry from the Anti-Phishing Working Group (APWG) \cite{apwg2023trends} documents millions of active phishing attacks annually, while the Verizon Data Breach Investigations Report (DBIR) \cite{verizon2024dbir} identifies stolen credentials as a primary initial breach vector across thousands of confirmed corporate security incidents.

Conventional defenses rely heavily on centralized signature databases and community blacklists (e.g., Google Safe Browsing, PhishTank). While blacklists provide high precision on cataloged domains, they exhibit significant discovery delays, often requiring hours to verify zero-hour domains \cite{sahoo2017malicious, rao2019detection, hannousse2021towards, marchal2017off}, during which zero-day attacks achieve maximum impact.

To overcome blacklist latency, machine learning (ML) has emerged as a proactive alternative, classifying domain, path, and structural indicators dynamically \cite{sahingoz2019machine, babagoli2019heuristic}. However, operational adoption in Security Operations Centers (SOCs) faces two critical bottlenecks: (i) \textbf{Model Opacity}, where opaque tree ensembles fail to provide interpretable feature attributions for containment decisions; and (ii) \textbf{Page Crawling Latency}, where full DOM rendering introduces substantial latency and risks drive-by malware \cite{opara2020htmlphish, adebowale2023intelligent}.

To address these challenges, this paper presents an empirical benchmark and explainable ML framework for phishing detection using 111 structured URL lexical, path, file, query, and resolver metrics on 58,645 URLs \cite{vrbancic2020datasets}. We evaluate five paradigms: Logistic Regression, Decision Trees, Random Forests, Linear SVM, and XGBoost across 11,729 held-out test samples with 5-fold cross-validation. To ensure transparency, we apply Tree-based Shapley Additive Explanations (TreeSHAP) \cite{lundberg2020local, lundberg2017unified} strictly on training data to ensure train-only attribution.

The primary contributions of this work are:
\begin{itemize}
    \item \textbf{Multi-Model Benchmark \& Statistical Validation}: We evaluate five supervised ML architectures across 11,729 test instances with 5-fold CV on training data ($N=46,916$), verifying that XGBoost's lead (5-seed test mean: $94.74\% \pm 0.08\%$ accuracy, $94.98\% \pm 0.08\%$ F1; seed 42: 94.82\% accuracy, 95.05\% F1; seeds vary model tree subsampling on a fixed split) over Random Forest is statistically significant via McNemar's test ($p < 0.001$).
    \item \textbf{Train-Only Attribution \& Ranking Stability}: We implement TreeSHAP strictly on training data to ensure train-only attribution, demonstrating robust ranking stability under model retraining ($\rho_s = 0.973 \pm 0.010$ across all 98 features, Kendall's $\tau = 0.795$ within the top 20, and 95.0\% top-20 set overlap across independent model training seeds; $\rho_s = 0.9991$ and $\tau = 0.9131$ under explanation subsampling) and identifying primary signals (directory length, domain activation age, delimiters).
    \item \textbf{Feature Reduction \& Cascaded Triage}: We evaluate nested subsets, selecting $k=20$ as the smallest tested subset within about 0.2\,pp of full training CV F1, which sustains full test performance (5-seed test mean: $94.71\% \pm 0.09\%$ accuracy, $94.95\% \pm 0.08\%$ F1) while lowering model-only latency to $1.67$\,ms per 1,000 queries. Furthermore, evaluating an uncertainty-gated cascade demonstrates that a zero-network lexical-only model (84 features) resolves 64.3\% of URLs in-memory with 97.78\% Tier-1 accuracy, supporting a tiered design for high-throughput perimeter triage that offloads 64.3\% of external resolver queries while maintaining 94.57\% overall pipeline accuracy.
    \item \textbf{Computational Profiling}: We measure model inference latency (mean $\pm$ std over 20 repeated runs), single-query latency ($N=1$), serialized model footprints, and throughput under batched execution.
\end{itemize}

The remainder of this manuscript is organized as follows: Section~\ref{sec:related_work} surveys related work. Section~\ref{sec:methodology} presents the methodology. Section~\ref{sec:experiments} details the setup. Section~\ref{sec:results} discusses results. Section~\ref{sec:discussion} addresses trade-offs and error analysis. Section~\ref{sec:conclusion} concludes the paper.

